\chapter{Reverse-Engineering the Targets}
\label{ch:targets}

The highest-leverage work in our entire campaign was done without fitting a
single model. It was \emph{forensic}: determining what the anonymized target
actually \emph{is}. This chapter teaches the toolkit through the definitive
worked example --- the community reverse-engineering of the Jane Street
responders (Kaggle discussion \#555562, John Payne, reproduced, verified, and
extended in our repository) --- and develops each tool with enough
mathematics that you can wield it on the next anonymized dataset you meet.

\section{Why forensics first}

An anonymized target is not an arbitrary label. Somebody \emph{computed} it
from real quantities, with real code, and then hid the code. Computation
leaves invariants: moving averages leave autocorrelation ramps; time shifts
leave lead--lag asymmetries; additive noise leaves spectral floors;
differences of related series leave cross-correlation fingerprints. Recovering
the construction tells you:

\begin{itemize}
\item \textbf{the ceiling} --- if the target is a forward-looking average of
  innovations, the innovations' unpredictability bounds every model
  (\S\ref{sec:ceilingderiv});
\item \textbf{the feature semantics} --- once you know the target is
  ``next-20-step average return,'' features cluster into interpretable roles
  (Chapter~\ref{ch:featforensics});
\item \textbf{legal transformations} --- deconvolution, timing alignment,
  and auxiliary-target construction (Chapter~\ref{ch:featurepatterns}) all
  require knowing the construction;
\item \textbf{what not to attempt} --- entire research programs (within-day
  target autoregression, for instance) are provably dead ends under the
  recovered structure.
\end{itemize}

\section{Tool 1: the autocorrelation function as a fingerprint reader}
\label{sec:acf}

\subsection{The mathematics of the fingerprint}

Let $s_t$ be a zero-mean white-noise sequence, $\Var(s_t)=\sigma^2$, and
define the $w$-step simple moving average
\[
r_t \;=\; \frac{1}{w}\sum_{j=0}^{w-1} s_{t+j}.
\]
The autocovariance of $r$ at lag $k \ge 0$ counts overlapping terms:
\begin{equation}
\Cov(r_t, r_{t+k})
= \frac{\sigma^2}{w^2}\,\max(0,\, w-k)
\qquad\Longrightarrow\qquad
\corr(r_t, r_{t+k}) = \max\!\Big(0,\ 1-\frac{k}{w}\Big).
\label{eq:smaacf}
\end{equation}
The autocorrelation of an SMA-of-white-noise is a \textbf{perfect triangular
ramp}: it decays \emph{linearly} and hits exactly zero at lag $w$. No AR
process, no trend, no seasonal pattern produces this shape. A linear ramp
with a sharp cutoff is as close to a smoking gun as time-series forensics
offers: \emph{this series is a moving average, and the cutoff lag is the
window length}.

Contrast the other canonical fingerprints you should have on file:

\begin{center}
\begin{tabularx}{\textwidth}{@{}lX@{}}
\toprule
ACF shape & generating process \\
\midrule
triangular ramp, hard zero at lag $w$ & SMA($w$) of white noise \\
geometric decay $\rho^k$ & AR(1) with coefficient $\rho$; ``level'' series \\
single spike at lag 0 & white noise (returns, innovations) \\
ramp that decays to a plateau $>0$ & SMA of noise \emph{plus} a persistent
level component \\
spikes at $k = P, 2P, \dots$ & seasonality with period $P$ \\
ramp with rounded (not sharp) corner & SMA plus additive observation noise
(the noise dilutes but does not move the corner) \\
\bottomrule
\end{tabularx}
\end{center}

\subsection{The worked example: three windows, three ramps}

Applied to the Jane Street responders (computed per symbol, within days,
averaged over thousands of symbol-days --- the averaging matters, single
series are too noisy):

\begin{itemize}
\item \resp{6}: linear decay, sharp cutoff at lag \textbf{20};
\item \resp{7}: same shape, cutoff at lag \textbf{120};
\item \resp{8}: same shape, cutoff at lag \textbf{4}.
\end{itemize}

By \eqref{eq:smaacf}: the three responders are moving averages of windows
20, 120, and 4. Moreover \resp{0}/\resp{3} share the lag-20 fingerprint,
\resp{1}/\resp{4} the lag-120 one, \resp{2}/\resp{5} the lag-4 one --- three
window lengths, three groups. Our own reproduction
(\code{scripts/profile\_features.py}, run over two 80-day windows) measures
ramp linearity as the $R^2$ of a straight-line fit to the pre-cutoff ACF:
\textbf{0.996} for \resp{6}, \textbf{0.99} for \resp{8}, with cutoffs at
19--20 and 4 respectively. The fingerprint is essentially perfect.

The host's own metadata corroborates: \code{responders.csv} tags
\{2,5,8\} with \code{tag\_1} (the SMA-4 group), \{0,3,6\} with \code{tag\_2}
(SMA-20), \{1,4,7\} with \code{tag\_3} (SMA-120), \{0,1,2\} with
\code{tag\_0} and \{3,4,5\} with \code{tag\_4}. Anonymized metadata often
confirms structure discovered numerically; read it \emph{after} forming
hypotheses, as a check against confirmation bias.

\subsection{Implementation: always via FFT}

Computing ACFs by looping over lags is $O(TK)$ per series and, worse,
invites off-by-one errors. Use the Wiener--Khinchin route: the
autocorrelation is the inverse FFT of the power spectrum.

\begin{lstlisting}
def mean_acf(X, max_lag):            # X: (n_series, T), rows demeaned/unit-std
    L  = 1 << int(np.ceil(np.log2(2 * X.shape[1])))   # zero-pad: no wraparound
    F  = np.fft.rfft(X, n=L, axis=1)
    ac = np.fft.irfft(np.abs(F)**2, n=L, axis=1)[:, :max_lag+1]
    ac /= ac[:, :1]                  # normalize by lag-0
    return ac.mean(axis=0)           # average the fingerprint across series
\end{lstlisting}

Zero-padding to at least $2T$ prevents circular wraparound; averaging across
(symbol, day) series turns a noisy per-series estimate into a clean
fingerprint. This exact function, applied to \emph{features} instead of
responders, produced one of our campaign's original discoveries
(Chapter~\ref{ch:featforensics}): three features carry the SMA fingerprint
with cutoff $\approx 160$.

\section{Tool 2: lead--lag regression --- discovering the time shift}
\label{sec:shift}

The second discovery in \#555562 is subtler and worth more. Payne regressed
the responders on the features at various relative shifts and found:
\textbf{features at time $t$ predict \resp{6} at time $t-20$ with linear
$R^2 \approx 0.5$} --- five hundred times better than they predict it at
$t+0$ (about $0.001$).

Think about what that asymmetry means. If the stored value at row $t$ were
the average of the \emph{past} 20 steps, contemporaneous features (which
describe the present and recent past) would correlate with it strongly at
lag 0. Instead the correlation peaks when the responder is taken from 20
rows \emph{earlier} --- that is, the value stored at row $t$ aligns with
information from rows $t$ through $t+20$. Conclusion: \textbf{the stored
responder is forward-shifted}. The value in column \resp{6}, row $t$, is
(approximately) the average of the underlying signal over $(t,\, t+20]$.
The competition's task, decoded: \emph{given features now, predict the
average of the signal over the next 20 steps.}

\begin{keybox}[: run the shift scan on every anonymized target]
For shifts $k \in \{-2w, \dots, +2w\}$, fit a fast linear model (ridge on a
feature subsample suffices) of $y_{t+k}$ on features$_t$ and plot $R^2(k)$.
A peak at $k=0$ means a nowcasting target; a peak at $k<0$ means the stored
target is forward-shifted (you are forecasting); the peak's offset measures
the shift. Ten minutes of compute; changes your entire mental model of the
problem.
\end{keybox}

Two corollaries of the discovered shift, both of which shaped our campaign:

\begin{enumerate}
\item \textbf{The features are mostly descriptions of the realized past.}
  If features explain the just-\emph{completed} window at $R^2=0.5$ but the
  upcoming window at $0.01$, then features are overwhelmingly
  backward-looking summaries --- ``nowcast'' information --- with only a
  sliver of genuine foresight. Feature forensics (Chapter
  \ref{ch:featforensics}) quantifies exactly which features carry the
  sliver.
\item \textbf{A high-SNR auxiliary task exists.} Predicting the realized
  window (the $R^2=0.5$ task) is legal as a \emph{training} objective ---
  it needs no future information at training time --- and provides
  gradient signal two orders of magnitude stronger than the real target.
  We wired this into our recurrent models as ``nowcast auxiliary heads''
  (Chapter~\ref{ch:featurepatterns}).
\end{enumerate}

\section{Tool 3: Fourier analysis --- confirming the construction}
\label{sec:fourier}

The spectral density of an SMA($w$) of white noise is
\[
S_r(f) \;\propto\; \left|\frac{\sin(\pi w f)}{w\,\sin(\pi f)}\right|^2,
\]
the squared Dirichlet kernel --- a $\mathrm{sinc}$-like envelope with nulls
at multiples of $1/w$. Payne fit exactly this envelope, nulls spaced $1/20$
apart, to \resp{6}'s spectrum. One mismatch: the observed troughs do not
reach zero. A pure SMA has \emph{exact} spectral nulls; a floor under the
nulls means \textbf{something broadband was added after the averaging} ---
observation noise. Simulating SMA-20 of Gaussian noise \emph{plus} a small
independent Gaussian term reproduced the observed spectrum almost perfectly.

So the full recovered construction is:
\begin{equation}
r^{(w)}_t \;=\; \frac{1}{w}\sum_{j=1}^{w} s_{t+j} \;+\; \eta_t,
\qquad w \in \{4, 20, 120\},
\label{eq:construction}
\end{equation}
with $s_t$ the underlying per-step signal (near-white, heavy-tailed ---
plausibly one-minute returns) and $\eta_t$ small additive noise (deliberate
obfuscation or measurement error). Two exchange variants exist ($A$: responders
6/7/8; $B$: responders 3/4/5, the same $s$ observed with independent venue
noise), and responders 0/1/2 are the $B{-}A$ differences at matching windows.

The additive noise term is not a footnote. It means: (i) even perfect
knowledge of $s$ cannot reach $\Rw = 1$ on $r$; (ii) any deconvolution of
$r$ back to $s$ (Chapter~\ref{ch:featforensics}) must be regularized, since
pure inversion amplifies $\eta$ at the spectral nulls; (iii) simulated
replicas of the data (Chapter~\ref{ch:synthetic}) must include it to be
faithful.

\section{The ceiling, derived}
\label{sec:ceilingderiv}

Now we can compute what Chapter~\ref{ch:problem} asserted. Under
\eqref{eq:construction}, the target at row $t$ averages \emph{future}
innovations $s_{t+1},\dots,s_{t+20}$. Condition on everything observable at
time $t$ (call it $\mathcal F_t$). Then
\[
\E\big[r^{(20)}_t \,\big|\, \mathcal F_t\big]
= \frac{1}{20}\sum_{j=1}^{20} \E\big[s_{t+j} \mid \mathcal F_t\big],
\]
and the achievable $R^2$ is
\[
R^2_{\max}
= \frac{\Var\!\big(\E[r_t \mid \mathcal F_t]\big)}{\Var(r_t)}
= \frac{\frac{1}{400}\Var\!\Big(\sum_j \E[s_{t+j}\mid\mathcal F_t]\Big)}
       {\frac{20\,\sigma_s^2}{400} + \sigma_\eta^2}.
\]
If innovations were exactly white ($\E[s_{t+j}\mid\mathcal F_t]=0$), the
ceiling is \emph{zero}. The observed leaderboard ceiling of $\approx 0.014$
measures precisely how non-white the innovations are: the market's per-step
returns have a faint predictable component --- from momentum/reversal
structure, cross-sectional spillovers, and features that genuinely lead ---
and \emph{that faint component, averaged over a 20-step window, is the
entire game}. Every point of leaderboard score is a statement about market
microstructure.

This derivation also kills a tempting research direction cleanly. Within a
day you never observe past values of $r$; but even if you did, under
\eqref{eq:construction} the target's own history is nearly useless:
$\Cov(r_t, r_{t+k}) = 0$ for $k \geq w$ by whiteness, so
$r_{t-20}$ tells you (almost) nothing about $r_t$. Payne said it directly
in the thread: even a perfect estimate of the responder twenty steps ago
has expectation zero for the current value. \emph{Do not build within-day
autoregressive models of this target.} We logged this as a
theorem-level negative result before spending any compute on it --- the
cheapest kind.

\section{Tool 4: cross-series relationships}

With per-series structure identified, examine relationships \emph{between}
series. The \#555562 analysis noticed, by overlaying lag-aligned plots, that
\resp{6} behaves like a ``derivative'' of \resp{7} (when the 20-average is
positive, the 120-average rises) and \resp{8} like a derivative of \resp{6}.
That is exactly what SMAs of one signal at nested windows do --- the shorter
average leads the longer one --- and it suggests immediately that
\textbf{differences of same-signal-different-window series are trend
signals}. Lagged by one day (the only legal access at inference), these
became our ``rsig'' feature block: $\text{SMA4}-\text{SMA20}$ (short
momentum), $\text{SMA20}-\text{SMA120}$ (medium momentum), $A{-}B$ at
matching window (venue spread). Measured value: $+0.0003$ on an XGBoost
harness where the nine raw lagged responders in un-differenced form
contributed \emph{nothing} ($-0.00003$). Construction knowledge converts
dead columns into live ones (Chapter~\ref{ch:featurepatterns} for the full
accounting).

\section{The forensic report template}

Every anonymized dataset deserves a written forensic report before modeling
begins. Ours (\code{docs/FEATURE\_RESEARCH.md} in the repository) follows a
template you can copy:

\begin{enumerate}
\item \textbf{Construction hypothesis} for each target: formula, window,
  shift, noise (equation \eqref{eq:construction} form), with the evidence
  (ACF cutoff, ramp linearity, spectral floor, shift-scan peak).
\item \textbf{Ceiling estimate} and its derivation.
\item \textbf{Consequences}: dead research directions (with proofs);
  legalized transformations (deconvolution, aux targets); alignment
  conventions every pipeline must respect.
\item \textbf{Stability checks}: every finding recomputed on a second,
  distant time window. (Our standard: two 80-day windows, 600 days
  apart. Structural facts --- windows, shifts --- reproduced exactly;
  that is what earns them load-bearing status.)
\item \textbf{Open questions}, ranked by expected value.
\end{enumerate}

\begin{drillbox}
The repository's \code{scripts/profile\_features.py} computes ACF ramps,
cutoffs, and linearity scores for all 88 columns in one run. Reproduce the
responder fingerprint table, then re-run on the first differences of the
responders. Predict what the ACF of a \emph{differenced} SMA-of-noise looks
like before you look --- deriving it from \eqref{eq:smaacf} is a five-line
exercise --- and check yourself. (You will meet the answer again in
Chapter~\ref{ch:smoothing}, where it explains why differencing pre-smoothed
features is the right anti-smoothing move.)
\end{drillbox}
